\documentclass[conference]{IEEEtran}
\usepackage{amsmath}
\usepackage{booktabs}
\usepackage{graphicx}

\title{Crypto3DStackCPU: A Secure Single-Layer CPU Stack with Sealed Images and AES Instruction Support}

\author{\IEEEauthorblockN{George David Tsitlauri}
\IEEEauthorblockA{
gdtsitlauri@gmail.com\\
gdtsitlauri.dev\\
github.com/gdtsitlauri
}}

\begin{document}
\maketitle

\begin{abstract}
Crypto3DStackCPU is a hardware/security research repository that implements a secure single-layer CPU stack with image sealing, authenticated headers, encrypted instruction fetch, encrypted-at-rest data handling, custom AES instructions, and contract-based validation. The repository includes a local assembler, end-to-end build and audit scripts, and HLS-oriented integration scaffolding. The strongest current evidence is implementation depth and security-flow completeness rather than broad benchmark coverage.
\end{abstract}

\section{Introduction}
Secure processor research often separates architecture, toolchain, and validation. Crypto3DStackCPU studies these concerns together in a compact repository that integrates assembly generation, image sealing, runtime validation, encrypted fetch, and negative tamper testing.

\section{Repository Scope}
The repository implements:
\begin{itemize}
\item a local assembler without MARS dependency,
\item sealed image construction with wrapped keys and authenticated headers,
\item encrypted text fetch and protected data handling,
\item custom AES instructions,
\item contract-driven execution validation,
\item a tamper-audit path that requires negative-test failure.
\end{itemize}

\section{Current Evidence}
The strongest artifacts are operational rather than benchmark-heavy. The repository includes:
\begin{itemize}
\item \texttt{run\_custom\_pipeline.ps1} for end-to-end assembly, encryption, and execution,
\item \texttt{run\_security\_audit.ps1} for positive and negative validation,
\item \texttt{3d\_test.cpp} for contract-based functional proof,
\item \texttt{3D\_hls\_component/} for HLS-facing integration work.
\end{itemize}

\section{Interpretation Boundary}
Crypto3DStackCPU should be framed as a serious implementation-heavy hardware/security extension repository. The current working tree does not yet provide a broad benchmark section across performance, area, or multiple workloads, so stronger comparative claims should wait for expanded evaluation.

\section{Conclusion}
Crypto3DStackCPU already demonstrates substantial hardware-security engineering. Its next step is not invention of core architecture, but systematic benchmarking and expanded manuscript-quality evaluation around the implemented secure execution path.

\bibliographystyle{IEEEtran}
\bibliography{references}
\end{document}
